% 天王星（综述）
% license CCBYSA3
% type Wiki

本文根据 CC-BY-SA 协议转载翻译自维基百科\href{https://en.wikipedia.org/wiki/Uranus}{相关文章}。
\begin{figure}[ht]
\centering
\includegraphics[width=8cm]{./figures/fb24dd1c06777ad6.png}
\caption{由旅行者 2 号拍摄的天王星真彩图像\(^\text{[a]}\)。其苍白而柔和的外观源于覆盖在云层之上的一层薄雾。} \label{fig_Uranus_1}
\end{figure}

天王星是距离太阳第七颗行星，它是一颗呈青绿色的气态冰巨星。其主要由处于超临界状态的水、氨和甲烷构成，在天文学中被称为“冰”或挥发物。该行星的大气层具有复杂的分层云结构，并在所有太阳系行星中拥有最低的最低温度（49 K（−224 °C，−371 °F））。天王星具有显著的 82.23° 自转轴倾角，并以 17 小时 14 分钟的周期作逆行自转。这意味着在其围绕太阳完成 84 个地球年的轨道周期中，其两极会各自经历约 42 年的连续日照，随后是约 42 年的连续黑夜。

在太阳系行星中，天王星的直径位居第三，质量位居第四。根据当前模型，在其富含挥发物的地幔层内部存在一个岩质核心，其外包覆着厚厚的氢—氦大气。在其高层大气中检测到痕量的碳氢化合物（被认为源于水解作用）以及一氧化碳和二氧化碳（被认为来自彗星）。天王星大气中存在许多尚未解释的气候现象，例如其高达 900 km/h（560 mph）的峰值风速、极冠的变化以及不规则的云层形成。该行星的内部热量也远低于其他巨行星，其成因至今仍不明。
